\section{Conclusions and outlook}\label{sec:concl}

\subsection{Summary of the main result}\label{sec:concl-thesis}

% code: solvers/contact/radial_chart_2d.py::NeuralRho2D
This work poses contact detection and computation as a boundary-to-boundary
\emph{transition map} on a learned coordinate atlas, and tests that map directly against the
ambient level set (a neural signed-distance function) on a single benchmark suite. The
governing idea is geometric rather than algorithmic. For two charted bodies $A$ and $B$, the
contact map composes the two boundary charts,
%
\begin{equation}\label{eq:concl-transition}
  \tmap = \map{B}^{-1}\circ\map{A}\,:\;\Theta_A \to \Theta_B,
\end{equation}
%
where $\map{A}$ is the chart of one boundary and $\map{B}^{-1}$ is the inverse chart of the
other, evaluated by Newton iteration. The gap and the contact normal then read directly from
$\tmap$, with the sign convention $\gn<0$ denoting penetration. Because the detector composes
two boundary charts, it never builds an ambient field and never crosses a medial axis. The
level set enters only as the comparison baseline; we do not advance it as a contribution of
this paper.

% code: common/models.py::MLP
The measurements support one conclusion: the \emph{geometric representation}, not the
constraint-enforcement scheme, controls fidelity on geometrically hard contact. A
fixed-capacity coordinate MLP that represents a body as an ambient field $\phi:\R^d\to\R$
learns low spatial frequencies first and high frequencies slowly or not at all
\citep{rahaman2019spectral,tancik2020fourier}. That field therefore smooths cusps, asperities,
and fine boundary detail, and it carries a non-smooth medial axis in any concavity. A boundary
chart instead stores the same body as a function \emph{one dimension lower} (a radius
$\rho:\Sph^{d-1}\to\R^+$, a height $h:\mathcal D\to\R$, or a boundary-fitted decoder
$D:\boldsymbol\xi\mapsto\xbf$), with no interior and no medial axis. Random-Fourier-feature
inputs remove the spectral-bias ceiling \citep{tancik2020fourier,mildenhall2020nerf}, so the
chart fits the $(d{-}1)$-dimensional boundary function sharply, exactly where the
$d$-dimensional field cannot.

% code: solvers/contact/chart_gap.py::evaluate_gap_chart
The quantitative ledger is consistent across the suite. On the convex, collinear, superposable
problems CV-1 through CV-4, the chart and the level set are equivalent and the chart buys
nothing: the chart-FEM contact solver reproduces Hertz line contact to $a(F)$--$E^*$ within
$1.6\%$ (CV-1), Cattaneo--Mindlin stick/slip to $5$--$11\%$ on a deep half-plane (CV-2), the
Brazilian-disc centre stress to $1.62\%$ (CV-3), and the nine-disc equibiaxial unit-cell stress
to $0.15\%$ (CV-4), matching but not beating the exact-SDF baseline
\citep{johnson1985contact,hondros1959evaluation,timoshenko1970elasticity,mindlin1949compliance}.
The two representations separate precisely where the geometry turns hard. On the cusped
non-convex superformula (CV-5), the neural SDF degrades to a gap error of $8\times10^{-3}L$ with an
out-of-plane-tilted normal, while the Fourier-feature radial chart reproduces the analytical
radial gap and normal to $3.8\times10^{-3}L$ and $0.42^\circ$, and drives the rigid-body
cam-drive dynamics to a $0.04\%$ trajectory match against the analytical chart
\citep{gielis2003generic}. On the self-similar Koch boundary (CV-6), a fixed-capacity SDF
($12{,}801$ parameters, trained on the exact Koch distance) plateaus at a capacity ceiling once
the boundary detail falls below its resolution, whereas the recursive IFS chart stores the
geometry in $O(1)$ and resolves any finite pre-fractal level exactly. At the true fractal limit
neither representation is well-posed, which we state rather than evade.

% code: solvers/fem/rough_block_decoder.py::RoughBlockDecoder
The capstone (CV-7) carries the comparison onto rough geometry in two complementary settings. On the
real Inada-granite profile, a self-affine tensile fracture \citep{inada2020digitalrocks} carried by a
one-dimensional Fourier height chart, the chart reconstructs the surface to $2.3~\mu$m where the ambient
neural SDF reaches only $107~\mu$m; the SDF smooths the mean asperity slope to $12.5^\circ$, and through
the Patton law that smoothing under-predicts the peak shear strength by $61\%$
\citep{patton1966multiple,barton1977shear}. On a synthetic band-limited surface carried by a
boundary-fitted Fourier-feature decoder, representative of the joint at the resolved scale, the decoder
reconstructs the geometry to $2.2\%$ RMS against $15.7\%$ for the SDF and $48\%$ for a plain tanh
decoder, and the chart-FEM solver on that decoder converges cleanly (frictionless contact Newton residual
$\sim10^{-9}$, MMS rates $\approx 2$, $\det J>0$). Under direct shear governed by \emph{plain Coulomb
friction} on the resolved faces, with no dilatancy law imposed, the apparent friction and the normal
dilation \emph{emerge} from the asperity geometry: the chart recovers a frictionless geometric
$\muapp\approx0.17$ and a frictional peak $\muapp\approx0.47$, while the SDF-smoothed surface yields
$\muapp\approx0.004$ and $0.31$, under-predicting dilatancy by $98\%$ and peak strength by $35\%$. The
two-deformable-block solve reproduces the
emergent strength to a mesh-converged peak $\muapp\approx0.53\pm2\%$ and the Patton anchor
$\tan(\phi_b{+}i)$ to $0.00\%$. These results substantiate the single mechanical claim of the
paper: the level set smooths the asperity slopes that set both strength and dilatancy, and the
transition-map chart resolves them.

\subsection{Outlook}\label{sec:concl-outlook}

% code: common/geometry.py::invert_decoder
Three concrete directions remain, each one removing a stated restriction. The present charts
assume star-shapedness (radial) or single-valuedness (height), which excludes overhangs and
re-entrant lobes. The route past this restriction is a multi-chart decoder atlas. A body is
covered by several boundary-fitted decoder charts $\{D_X\}$, and their overlaps reconcile
through Newton-inverted transition maps of the form \eqref{eq:concl-transition}: a query handled
by chart $X$ on one body carries into chart $Y$ on the other through $\map{Y}^{-1}\circ\map{X}$,
even where no single chart is global. The supporting components already exist in the framework:
the decoder inversion of \code{invert\_decoder}, the SVD-stabilized Jacobian operators, and a
persistent-homology chart-spawn pipeline that opens chart pairs from $H_0$ contact events
\citep{cohensteiner2007stability,liu2020implicit}. We do not exercise them in the present
results, so the extension is largely one of orchestration rather than of new theory. It
lifts the single-valuedness restriction that confines the rock-joint study to non-overhanging
surfaces.

% code: benchmarks/contact/cv_numerical/cv3_brazilian_fem.py::run_cv1
A second direction is full three-dimensional axisymmetric Hertz contact with local mesh
refinement at the contact patch. Patch resolution, not the formulation, limited the
two-dimensional line-contact verifications (CV-1, CV-2); an adaptively refined axisymmetric
solve would close that gap and place the chart-FEM contact pressure against the closed-form
$p_0$ and $a(F)$ in 3-D \citep{johnson1985contact,kikuchi1988contact}. The same refinement
strategy applies to the real-contact-area estimate on the rough joint, where the present
node-based active-set fraction $A_c/A$ is mesh-sensitive even though its trend is robust:
partial contact localizes onto fewer asperity ridges under shear.

% code: benchmarks/contact/cv_numerical/rock_joint_decoder_cyclic.py::run_cyclic
A third direction is to harden the finite-deformation and cyclic/dynamic extensions reported as
supplementary phases. The cyclic shear already produces stick/slip hysteresis with a per-cycle
energy ledger that closes to $\sim1\%$ and Plesha-type asperity-degradation decay
\citep{plesha1987constitutive,goodman1976methods}. The two-deformable moving-master
node-to-surface Coulomb contact, however, leaves a $1.5$--$6\%$ friction residual that does not
shrink under refinement; this residual is a non-smoothness artifact, not a discretization error. A
semismooth-Newton or augmented-Lagrangian treatment with mortar coupling
\citep{alart1991mixed,simo1992augmented,puso2004mortar,laursen2002computational} is the path to
tight convergence under cyclic and frictional loading. The explicit chart-MPM cross-check
reproduces the Coulomb floor ($\muapp\approx0.34$ against a base $\mu=0.4$) but not dilatancy,
because a mated rough-on-rough block is unstable in explicit penalty dynamics and requires
implicit or mortar contact. Embedding the transition-map detector \eqref{eq:concl-transition}
inside such a stabilized dynamic solver, together with a finite-strain hyperelastic interior
\citep{ball1977convexity,ciarlet1988mathematical}, would extend the
representation-controls-fidelity result from quasi-static rock joints to large-sliding,
finite-deformation contact. Across all three directions the conclusion stays constant: fidelity on
hard contact follows the geometry one chooses to resolve, and the transition map resolves what
the level set smooths.
